% Generated draft: requires mathematical and source review.
\documentclass[11pt]{article}
\usepackage[margin=25mm]{geometry}
\usepackage{fontspec}
\setmainfont{Latin Modern Roman}
\usepackage{amsmath,amssymb,mathtools,unicode-math}
\setmathfont{Latin Modern Math}
\usepackage{enumitem,needspace,xurl,hyperref}
\setlength{\parindent}{0pt}
\setlength{\parskip}{5pt}
\setlength{\emergencystretch}{3em}
\newcommand{\N}{\mathbb{N}}
\newcommand{\Z}{\mathbb{Z}}
\newcommand{\Q}{\mathbb{Q}}
\newcommand{\R}{\mathbb{R}}
\newcommand{\C}{\mathbb{C}}
\newcommand{\F}{\mathbb{F}}
\newcommand{\A}{\mathbb{A}}
\newcommand{\PP}{\mathbb{P}}
\newcommand{\E}{\mathbb{E}}
\newcommand{\eps}{\varepsilon}
\newcommand{\dd}{\,\mathrm{d}}
\begin{document}
\section*{UnsolvedMath ID: 30000337}
\textbf{Source problem code:} catalogue entry
\subsection{Definitions and mathematical statement}
Let $S$ be the surface in the source and let $G(S)$ be its intended connected graph of geometric objects, with the unit-edge path metric $d_G$ and mapping-class-group action. For an essential subsurface $X$, write $\pi_X$ for the prescribed subsurface projection into its curve/arc projection graph $\mathcal C(X)$. Call $X$ a hole when $\pi_X(v)\ne\varnothing$ for every vertex $v\in G(S)$. With the usual bounded-set projection convention, put
\[
 d_X(\alpha,\beta)=\operatorname{diam}_{\mathcal C(X)}
                         (\pi_X(\alpha)\cup\pi_X(\beta)).
\]
Annular projection targets and the allowed surface types must follow the source construction.

The first question is whether pairwise intersection of holes forces $G(S)$ to be Gromov hyperbolic, meaning that its geodesic triangles are $\delta$-thin for some uniform $\delta$. The exact meaning of intersection must be fixed: essential intersection or overlap of subsurfaces, nesting and choice of isotopy representatives can lead to different conditions.

The second question asks for a distance formula. For a threshold $A>0$, define $[r]_A=r$ when $r\ge A$ and zero otherwise. The desired uniform comparison is that suitable constants $A,K,C$, with $K\ge1$, $C\ge0$, give
\[
 K^{-1}\sum_{X\text{ a hole}}[d_X(\alpha,\beta)]_A-C
 \le d_G(\alpha,\beta)
 \le K\sum_{X\text{ a hole}}[d_X(\alpha,\beta)]_A+C
\]
for every pair of vertices. The sum must be well defined and finite, or otherwise appropriately convergent, for each pair. Constants may depend on the surface and specified graph class, not on $\alpha,\beta$.

The extract does not define the admissible graph families, vertex and edge types, coarse-projection axioms, hole equivalence, or whether both questions share precisely the same hypotheses. A mapping-class-group action and an unspecified projection alone do not identify one graph-theoretic theorem. The two displayed targets retain the source programme without assuming that every arbitrary graph with such an action satisfies a distance formula or is hyperbolic.

The literature assessment checked on 21 August 2026 reports Masur--Minsky-style positive results for some complexes but says the disk-graph extension in the passage remains difficult and unresolved. It supplies no complete graph class or modern universal theorem. This is dated triage requiring independent verification.
\subsection{Short English statement}
For the intended surface graphs, does pairwise intersection of projection holes imply hyperbolicity, and do large projection distances over all holes determine graph distance up to uniform constants? The graph and projection hypotheses must be recovered rather than inferred from the group action alone.
\subsection{Sources}
\begin{itemize}
\item \emph{Low-Dimensional Manifolds} (2006), Oberwolfach Reports, DOI 10.4171/owr/2005/45, pp. 2519--2570, \url{https://doi.org/10.4171/owr/2005/45}; proposed year supplied as 2005.
\item The sole literature source in the assessment checked 21 August 2026 is the original report, \url{https://doi.org/10.4171/owr/2005/45}. It mentions Masur--Minsky-style results and a disk-graph extension without full additional references.
\item Local UnsolvedMath registry, permanent ID $30000337$, exact source code \texttt{catalogue entry}, created 25 August 2026 and updated 16 September 2026. Corrected-and-verified status records meaning-preserving editorial and notation cleanup with both questions retained.
\end{itemize}
\end{document}
